\chapter{Результаты}
\label{ch:results}

Для формирования и определения радиусов колец Ньютона была использована установка, состоящая из двух линз: конденсатора и объектива, призмы прямого света и измерительного микроскопа с полупрозрачной пластинкой и линзой, радиус которой был определен в данной работе. В качестве источника исследуемого света использовалась ртутная лампа. В микроскопе присутствует окулярная шкала для измерения размеров наблюдаемых интерференционных картин. Подробное описание установки и калибровки шкалы см. в Приложении.

% Цена деления окулярной шкалы микроскопа определилась при помощи совмещения ее с калиброванной объектной шкалой. Полный метод калибровки см. в Приложении.

В результате были получены радиусы темных и светлых колец Ньютона для ртутной линии с длиной волны $546,1 \ \text{нм} ^{\text{\cite{Lab}}}$ в зависимости от их номера. Данные значения представлены на рис. \ref{fig:r2_m}.

\begin{figure}[ht]
    \begin{gnuplot}[terminal=pdf]
        set xl "m, номер кольца"
        set xr [0:13]

        set yl "r^2, мм^2"
        set yr [0:0.1]
        set format y "%.2f"

        set key b r
        set grid

        del = 0.98*(10**(-6))

        rd(x) = kd*x + bd
        fit rd(x) "txt/r(m).txt" u 1:(($2*del)**2*(10**6)) via kd,bd
        rw(x) = kw*x + bw
        fit rw(x) "txt/r(m).txt" u 1:(($3*del)**2*(10**6)) every ::1 via kw,bw

        plot rd(x)  lw 2 lc "#7000ff"  t "темные кольца", rw(x)  lw 2 lc "#ffc000"  t "светлые кольца", \
            "txt/r(m).txt" u 1:(($2*del)**2*(10**6)):((5/$2)*($2*del)**2*(10**6)) w ye  pt 0 lw 1 lc "#c000ff"  not, \
            "txt/r(m).txt" u 1:(($3*del)**2*(10**6)):((5/$3)*($3*del)**2*(10**6)) every ::1 w ye  pt 0 lw 1 lc "#ff8000"  not
    \end{gnuplot}
    \vspace{-30pt}
    \caption{Квадраты радиусов $r$ темных и светлых колец зеленого цвета в зависимости от их номера $m$.}
    \vspace{-10pt}
    \label{fig:r2_m}
\end{figure}
\FloatBarrier

По графику было определено, что зависимости имеют линейный вид. Из этого был сделан вывод о постоянности отношения радиуса колец к их номеру. Также было замечено, что наклон графиков для темных и светлых колец схож, поэтому с помощью МНК был определен их средний коэффициент $k \approx (7.135 \pm 0.021) \ 10^{-3} \, \text{мм}^2$.
Благодаря этому значению было получено, что радиус линзы микроскопа равен $(12,98 \pm 0,05) \ \text{мм}$. Это показало, что изучение колец Ньютона позволяет бесконтактно определять радиусы линз с высокой точностью.

% $k_\text{темн} = (7.135 \pm 0.021) \ 10^{-3} \, \text{мм}^2$ \\
% $k_\text{свет} = (7.04 \pm 0.04) \ 10^{-3} \, \text{мм}^2$ \\
% $k = (7.09 \pm 0.03) \ 10^{-3} \, \text{мм}^2$ \\

Также были проведены измерения ''биений'' между ртутными линиями с длинами волн ($546,1 \ \text{нм} ^{\text{\cite{Lab}}}$) и ($579,1 \ \text{нм} ^{\text{\cite{Lab}}}$). Из этого была получена разность радиусов двух соседних четких систем колец. С помощью данного значения была рассчитана разница длин волн $\Delta \lambda$ между спектральными компонентами, равная $(34,1 \pm 2,1) \ \text{нм}$. Это значение сходится в пределах погрешности с табличным $\Delta \lambda = 579,1 - 546,1 = 33,0 \ \text{нм}$.
Это показало, что данный метод можно использовать для определения разности длин волн двух спектральных компонент источника света.

% $\Delta m = 16 \pm 1$, $\lambda_2 = 546,1 \ \text{нм}$ - зеленый цвет \cite{Lab} \\
% $\Delta \lambda = \lambda_2 / \Delta m \approx (34,1 \pm 2,1) \ \text{нм}$ \\
% $\lambda_1 = (580,2 \pm 2,1) \ \text{нм}$ \\
% ближайшее табличное значение \cite{Lab}: $\lambda_1 = 579,1 \ \text{нм}$ - желтый цвет \\

% $R = k / \lambda_2 \approx (12,98 \pm 0,05) \ \text{мм}$
